\documentclass[11pt]{article}

% ===================== PACKAGES =====================
\usepackage[a4paper,margin=1in]{geometry}
\usepackage{amsmath,amssymb}
\usepackage{enumitem}
\usepackage{fancyhdr}
\usepackage{xcolor}
\usepackage{booktabs}
\usepackage{array}
\usepackage{longtable}
\usepackage{tabularx}

% ===================== HEADER & FOOTER =====================
\pagestyle{fancy}
\fancyhf{}
\lhead{CSE400: Fundamentals of Probability in Computing}
\rhead{Milestone 4 Scribe}
\cfoot{\thepage}

% ===================== CUSTOM COMMANDS =====================
\newcommand{\scribequestion}[1]{
    \vspace{0.5cm}
    \noindent\textbf{\large #1}
    \vspace{0.2cm}
}

\newcommand{\subquestion}[1]{
    \vspace{0.3cm}
    \noindent\textbf{#1}
    \vspace{0.1cm}
}

\newcommand{\highlight}[1]{\textit{#1}}

% ===================== TITLE =====================
\title{
    \normalsize School of Engineering and Applied Science (SEAS), Ahmedabad University \\
    \vspace{0.2cm}
    \textbf{CSE400: Fundamentals of Probability in Computing}\\
    \Large Milestone 4 Scribe\\
    \large \textit{Handover Probability in Drone Cellular Networks}
}
\author{}
\date{}

\begin{document}
\maketitle

\vspace{-2cm}
\begin{center}
    \begin{tabular}{ll}
        \textbf{Group:} & {s1 g14 net\hspace{3in}}
    \end{tabular}
\end{center}

\hrule
\vspace{0.5cm}

\begin{enumerate}[label=\textbf{Scribe Question \arabic*:}, leftmargin=*, itemsep=1.5em]

% ---------------- SCRIBE QUESTION 1 ----------------
\item \textbf{Baseline System and Motivation for Randomization}

\subquestion{What was the deterministic baseline?}

In Milestone 3, we built a simulation of a drone cellular network where UAVs act as flying base
stations serving a single stationary user on the ground. The drones' starting positions were placed
using a Poisson Point Process at a fixed density, and each drone then travels in a straight line at a
random angle.

The system always connects the user to whatever drone happens to be physically closest at that
moment --- this is the nearest-neighbor association policy:
\[
x^*(t) = \arg\min_i \|x_i(t)\|
\]
A handover gets recorded every time the identity of that closest drone changes. We looked at two
movement scenarios: SSM, where every drone flies at the same fixed speed, and DSM, where each
drone is assigned its own speed sampled from either a Rayleigh or Uniform distribution. Every
parameter in M3 was held constant, matching the values confirmed in the paper and slides:

\vspace{0.3cm}
\begin{center}
\begin{tabular}{ll}
    \toprule
    \textbf{Parameter} & \textbf{M3 Fixed Value} \\
    \midrule
    Drone density $\lambda_0$ & $1.0$ DBS/km$^2$ $= 1 \times 10^{-6}$ /m$^2$ \\
    Speed (SSM) & $45$ km/h ($12.5$ m/s) \\
    Exclusion radius $u_0$ & $500$ m \\
    Time window $T$ & $400$ s \\
    Monte Carlo runs & $1{,}00{,}000$ simulations \\
    \bottomrule
\end{tabular}
\end{center}
\vspace{0.3cm}

\subquestion{Why introduce randomization?}

Running everything at one fixed set of values only tells us the theory holds for that single scenario.
In practice, drone networks look different every time they're deployed --- the number of drones, their
speeds, and the distances involved all vary. When we instead draw parameters randomly each run,
every simulation becomes its own independent test. If the theoretical predictions still hold up across
all of those varied runs, we can be confident the results genuinely generalize rather than being a
coincidence of one carefully chosen configuration.

% ---------------- SCRIBE QUESTION 2 ----------------
\newpage
\item \textbf{Randomized Algorithm Design}

\subquestion{What component now involves randomness?}

We built a new central module called \textbf{randomised\_config.py} that executes at the very
beginning of each simulation. Rather than loading hard-coded constants, it samples five key physical
parameters fresh from Uniform distributions every run. Everything else --- the nearest-neighbor
connection rule, the handover detection logic, and how the drones move --- stays exactly the same
as it was in M3.

\subquestion{What distributions are used?}

\vspace{0.3cm}
\begin{center}
\begin{tabular}{lll}
    \toprule
    \textbf{Parameter} & \textbf{Distribution} & \textbf{Range (from code)} \\
    \midrule
    Drone density $\lambda_0$ & Uniform & $5\times10^{-7}$ to $2\times10^{-6}$ UAVs/m$^2$ \\
    Mean speed $v^{\text{Mean}}$ & Uniform & $30$ to $60$ km/h \\
    Speed spread $v^{\text{Diff}}$ & Uniform & $20$ to $120$ km/h \\
    Exclusion radius $u_0$ & Uniform & $300$ to $800$ m \\
    Time window $T$ & Uniform (integer) & $200$ to $400$ s \\
    \bottomrule
\end{tabular}
\end{center}
\vspace{0.3cm}

The specific values that were actually drawn for the run used to produce our graphs ---
verified against the saved JSON file, were:

\vspace{0.3cm}
\begin{center}
\begin{tabular}{lll}
    \toprule
    \textbf{Parameter} & \textbf{M3 Deterministic (Fixed)} & \textbf{M4 Randomized (Actual Draw)} \\
    \midrule
    $\lambda_0$ & $1.00 \times 10^{-6}$ /m$^2$ & $1.44 \times 10^{-6}$ /m$^2$ \\
    Mean speed   & $45.0$ km/h & $53.1$ km/h \\
    $u_0$        & $500$ m & $335.8$ m \\
    $T$          & $400$ s & $272$ s \\
    Simulations  & $1{,}00{,}000$ & $10{,}000$ \\
    \bottomrule
\end{tabular}
\end{center}
\vspace{0.3cm}

\subquestion{How does it differ from the deterministic baseline?}

In M3, the exact same numbers fed into every single run. In M4, each run is genuinely a fresh
experiment with its own parameter set. To keep the comparison fair, a \textbf{save-and-load system}
is used: the sampled values are written to a JSON file immediately, and every theory script then reads
from that same file, so the simulation and the theoretical calculations are always working from
identical inputs.

% ---------------- SCRIBE QUESTION 3 ----------------
\newpage
\item \textbf{Probabilistic Reasoning Behind the Algorithm}

\subquestion{How does randomness help address uncertainty?}

In a real-world drone network, nothing stays fixed --- drone counts, speeds, and deployment distances
all shift from one situation to the next. If we only ever test at one specific configuration, the most we
can conclude is that the theory works there. By sampling parameters from distributions, each run
probes the theory at a completely different operating point. Agreement across all of those varied
conditions is a much stronger form of validation.

\subquestion{What probabilistic models support this design?}

\begin{itemize}[leftmargin=*, itemsep=0.3em]
    \item \textbf{PPP for drone deployment:} The paper's derivations hold for any density $\lambda_0$,
          so randomizing it simply tests the same theory at a different point rather than violating
          any assumption.

    \item \textbf{Displacement Theorem (Lemma 1):} This result guarantees that drones remain
          Poisson-distributed at every instant regardless of how they move, for any $\lambda_0$ ---
          our randomized value still falls within that guarantee.

    \item \textbf{Speed PDF for DSM (Lemma 3):} Lemma 3 applies to any valid speed distribution
          $f_v(v)$, so choosing a different Rayleigh sigma $\sigma$ just plugs a different valid
          distribution into the same formula without breaking anything.

    \item \textbf{Exclusion radius $u_0$:} Varying $u_0$ across the $300$--$800$ m range captures
          the reality that drones don't always land at the same distance from a user.
\end{itemize}

\subquestion{Why is this expected to improve the analysis?}

The randomized approach doesn't change the underlying handover physics. What it does is make our
conclusions far more defensible. A finding that holds across many independently drawn parameter
settings is substantially more convincing than one that was only ever verified at a single
hand-picked configuration.

